\documentclass[11pt]{article}

\usepackage[margin=1in]{geometry}
\usepackage{amsmath}
\usepackage{booktabs}
\usepackage{microtype}
\usepackage[hidelinks]{hyperref}
\usepackage{xcolor}

\definecolor{cmxred}{RGB}{210,25,20}

\title{Where the Error Bars on the Red CMX Curve Come From}
\author{}
\date{}

\begin{document}
\maketitle

\section*{Short answer}

In the LiH absolute-energy-error plot, the \textcolor{cmxred}{red circles and
dashed curve} are energies obtained after connected moments expansion (CMX)
error mitigation.  Their vertical error bars represent uncertainty caused by
finite statistical sampling of the quantum measurements.  They are not the
difference between the red point and the exact ground-state energy.

\section*{Quantities measured for CMX}

CMX estimates the ground-state energy from three Hamiltonian moments:
\begin{equation}
  a=\langle \hat H\rangle,\qquad
  b=\langle \hat H^2\rangle,\qquad
  c=\langle \hat H^3\rangle .
\end{equation}
Each expectation value is estimated from a finite number of measurement shots.
Consequently, $a$, $b$, and $c$ have statistical variances
$\operatorname{Var}(a)$, $\operatorname{Var}(b)$, and
$\operatorname{Var}(c)$.

\section*{Propagation into the CMX energy}

The CMX energy is a nonlinear function
\begin{equation}
  E_{\mathrm{CMX}}=f(a,b,c).
\end{equation}
The paper propagates the variances of the measured moments through this
function:
\begin{equation}
\label{eq:propagation}
  \operatorname{Var}(E_{\mathrm{CMX}})
  =
  \left(\frac{\partial E_{\mathrm{CMX}}}{\partial a}\right)^2
       \operatorname{Var}(a)
  +
  \left(\frac{\partial E_{\mathrm{CMX}}}{\partial b}\right)^2
       \operatorname{Var}(b)
  +
  \left(\frac{\partial E_{\mathrm{CMX}}}{\partial c}\right)^2
       \operatorname{Var}(c).
\end{equation}
The plotted error-bar magnitude is therefore the propagated standard
uncertainty
\begin{equation}
  \boxed{\sigma_{E_{\mathrm{CMX}}}
  =\sqrt{\operatorname{Var}(E_{\mathrm{CMX}})}}.
\end{equation}
Thus, a red point at energy $E_{\mathrm{CMX}}$ is displayed with a vertical
range conventionally interpreted as
$E_{\mathrm{CMX}}\pm\sigma_{E_{\mathrm{CMX}}}$.

Equation~\eqref{eq:propagation} is Eq.~(72) in the paper's Supplementary
Information.  Its expression specialized to the second-order CMX estimator is
given in Supplementary Eq.~(73).  The authors state that the CMX error bar is
calculated using Eq.~(73).  Because this propagation omits covariance terms,
they also note that the actual variance could be smaller.

\section*{Values for the LiH red curve}

The paper's \emph{Source Data Fig.~3} spreadsheet labels these numbers
``REM+CF+CMX error bar.''  For the LiH panel, their magnitudes are:

\begin{center}
\begin{tabular}{cc}
\toprule
Bond distance (\AA) & Error-bar magnitude (mHa) \\
\midrule
0.5 & 1.747 \\
0.7 & 4.135 \\
0.9 & 1.496 \\
1.2 & 5.452 \\
1.5 & 1.525 \\
1.8 & 1.711 \\
2.2 & 1.615 \\
2.6 & 2.280 \\
3.0 & 1.843 \\
\bottomrule
\end{tabular}
\end{center}

Here $1~\mathrm{mHa}=10^{-3}$ Hartree.  These values quantify statistical
precision.  By contrast, the vertical coordinate in the absolute-error panel
is
\begin{equation}
  \left|E_{\mathrm{CMX}}-E_{\mathrm{exact}}\right|,
\end{equation}
which quantifies agreement with exact diagonalization.  These are two
different notions of error.

\section*{Source}

S.~Guo \emph{et al.}, ``Experimental quantum computational chemistry with
optimised unitary coupled cluster ansatz,'' \emph{Nature Physics} \textbf{20},
1240--1246 (2024),
\href{https://doi.org/10.1038/s41567-024-02530-z}
{doi:10.1038/s41567-024-02530-z}.  See Supplementary Eqs.~(72)--(73) and
Source Data Fig.~3.

\end{document}
